\documentclass[10pt,a4paper]{amsart}
\usepackage{amsmath,amssymb,amsthm,graphicx}
\usepackage{verbatim}
\newenvironment{reference}{\comment}{\endcomment}
%\newenvironment{reference}{}{}
\newenvironment{slogan}{\comment}{\endcomment}
\newenvironment{history}{\comment}{\endcomment}

% For commutative diagrams we use Xy-pic
\usepackage[all]{xy}

% We use 2cell for 2-commutative diagrams.
\xyoption{2cell}
% Theorem environments.
%
\theoremstyle{plain}
\newtheorem{theorem}[subsection]{Theorem}
\newtheorem{proposition}[subsection]{Proposition}
\newtheorem{lemma}[subsection]{Lemma}

\theoremstyle{definition}
\newtheorem{definition}[subsection]{Definition}
\newtheorem{example}[subsection]{Example}
\newtheorem{exercise}[subsection]{Exercise}
\newtheorem{situation}[subsection]{Situation}

\theoremstyle{remark}
\newtheorem{remark}[subsection]{Remark}
\newtheorem{remarks}[subsection]{Remarks}

\numberwithin{equation}{subsection}

% Macros
%
\def\lim{\mathop{\mathrm{lim}}\nolimits}
\def\colim{\mathop{\mathrm{colim}}\nolimits}
\def\Spec{\mathop{\mathrm{Spec}}}
\def\Hom{\mathop{\mathrm{Hom}}\nolimits}
\def\Ext{\mathop{\mathrm{Ext}}\nolimits}
\def\SheafHom{\mathop{\mathcal{H}\!\mathit{om}}\nolimits}
\def\SheafExt{\mathop{\mathcal{E}\!\mathit{xt}}\nolimits}
\def\Sch{\mathit{Sch}}
\def\Mor{\mathop{\mathrm{Mor}}\nolimits}
\def\Ob{\mathop{\mathrm{Ob}}\nolimits}
\def\Sh{\mathop{\mathit{Sh}}\nolimits}
\def\NL{\mathop{N\!L}\nolimits}
\def\CH{\mathop{\mathrm{CH}}\nolimits}
\def\proetale{{pro\text{-}\acute{e}tale}}
\def\etale{{\acute{e}tale}}
\def\QCoh{\mathit{QCoh}}
\def\Ker{\mathop{\mathrm{Ker}}}
\def\Im{\mathop{\mathrm{Im}}}
\def\Coker{\mathop{\mathrm{Coker}}}
\def\Coim{\mathop{\mathrm{Coim}}}

% Boxtimes
%
\DeclareMathSymbol{\boxtimes}{\mathbin}{AMSa}{"02}

%
% Macros for moduli stacks/spaces
%
\def\QCohstack{\mathcal{QC}\!\mathit{oh}}
\def\Cohstack{\mathcal{C}\!\mathit{oh}}
\def\Spacesstack{\mathcal{S}\!\mathit{paces}}
\def\Quotfunctor{\mathrm{Quot}}
\def\Hilbfunctor{\mathrm{Hilb}}
\def\Curvesstack{\mathcal{C}\!\mathit{urves}}
\def\Polarizedstack{\mathcal{P}\!\mathit{olarized}}
\def\Complexesstack{\mathcal{C}\!\mathit{omplexes}}
% \Pic is the operator that assigns to X its picard group, usage \Pic(X)
% \Picardstack_{X/B} denotes the Picard stack of X over B
% \Picardfunctor_{X/B} denotes the Picard functor of X over B
\def\Pic{\mathop{\mathrm{Pic}}\nolimits}
\def\Picardstack{\mathcal{P}\!\mathit{ic}}
\def\Picardfunctor{\mathrm{Pic}}
\def\Deformationcategory{\mathcal{D}\!\mathit{ef}}
\usepackage{cleveref}
\crefname{originallemma}{original lemma}{original lemmas}
\Crefname{originallemma}{Original lemma}{Original lemmas}
\makeatletter
\newlabel{lemma-Zariski-topology}{{00E0}{source}{}{sourcepointer-1}{}}
\newlabel{lemma-Zariski-topology@cref}{{[originallemma][0][]00E0}{[1][0][]source}{}{}{}}
\makeatother
\crefname{originaldefinition}{original definition}{original definitions}
\Crefname{originaldefinition}{Original definition}{Original definitions}
\makeatletter
\newlabel{topology-definition-space-jacobson}{{005U}{source}{}{sourcepointer-2}{}}
\newlabel{topology-definition-space-jacobson@cref}{{[originaldefinition][0][]005U}{[1][0][]source}{}{}{}}
\makeatother
\title{A ring is Jacobson exactly when its spectrum is a Jacobson topological space}
\author{The Stacks Project Authors}
\date{stacks-project:algebra.tex:local-snapshot-20261012; GNU FDL 1.2 or later}
\begin{document}
\maketitle
\noindent\textit{Editorial note.} All delivered mathematical text is the author's own, copied from the pinned Stacks Project chapter file algebra.tex; no statement or proof step was rewritten, reordered or removed.\par\medskip
\noindent\textit{Editorial note.} Preamble provenance, as required for this source: the frozen snapshot stages the chapter's own preamble file beside it (preamble.tex, 217 lines, 8116 bytes, the same shared Stacks-wide preamble the published chow and etale-cohomology cards use); the first line of the pinned chapter file is \textbackslash{}input\{preamble\}, so no preamble text lives inside the chapter itself. The card takes that file's lines 8-18 (the verbatim package and the three \textbackslash{}comment-based environments reference/slogan/history at lines 9-12, section 13.2, together with the Xy-pic setup) and its lines 154 to the runtime end of the file (the native theorem declarations and mathematical macros). Not one character was changed, no package was added, no macro was defined, and the excerpt serves only as preamble, never as a span coordinate.\par\medskip
\noindent\textit{Editorial note.} The same-chapter pointers are located at their \textbackslash{}label lines inside the pinned chapter file; the cross-chapter pointers are located in the project's own tag list staged in the snapshot (sources/stacks-chow/tags.txt), at the verbatim lines carrying their canonical tags.\par\medskip
\noindent\textit{Editorial note.} The counted statement and proof refer to a definition of the Topology chapter and to a same-chapter lemma; both are recorded as navigation stubs pointed at their \textbackslash{}label lines, the cross-chapter one carrying kind definition.\par\medskip
\noindent\textit{Editorial note.} The counted spans contain no \textbackslash{}cite, so the card carries no bibliography.\par\medskip
\noindent\textit{Allowed standard prerequisites (editorial).}\begin{itemize}
\item The topological notion of a Jacobson space. Invoked in the statement and used at the start of the counted proof (lines 6910-6919). Supplies the exact condition on closed subsets of the spectrum that has to be matched.
\item Zariski closed subsets correspond to radical ideals. Invoked at the end of the counted proof (lines 6939-6941). Converts the equality of two closed subsets into equality of the corresponding radical ideals.
\end{itemize}\medskip
\ifdefined\hypertarget\hypertarget{sourcepointer-1}{}\fi
\noindent\textit{Source reference.} Original lemma 00E0: sources/stacks-algebra/pinned-original-source.tex, source lines 3204--3204. Navigation pointer for \textbackslash{}ref\{lemma-Zariski-topology\} used inside the counted span (line 6924); the argument that consumes it is retained in the span.\par\medskip
\ifdefined\hypertarget\hypertarget{sourcepointer-2}{}\fi
\noindent\textit{Source reference.} Original definition 005U: sources/stacks-chow/tags.txt, source lines 222--222. Navigation pointer for \textbackslash{}ref\{topology-definition-space-jacobson\} used inside the counted span (line 6912); the argument that consumes it is retained in the span.\par\medskip
\begin{lemma}
\label{lemma-jacobson}
A ring $R$ is Jacobson if and only if $\Spec(R)$
is Jacobson, see Topology,
Definition \ref{topology-definition-space-jacobson}.
\end{lemma}

\begin{proof}
Suppose $R$ is Jacobson. Let $Z \subset \Spec(R)$
be a closed subset. We have to show that the set of closed
points in $Z$ is dense in $Z$. Let $U \subset \Spec(R)$
be an open such that $U \cap Z$ is nonempty.
We have to show $Z \cap U$ contains a closed point
of $\Spec(R)$. We may
assume $U = D(f)$ as standard opens form a basis for the
topology on $\Spec(R)$. According to
Lemma \ref{lemma-Zariski-topology} we may assume that
$Z = V(I)$, where $I$ is a radical ideal. We see also
that $f \not \in I$. By assumption, there exists a
maximal ideal $\mathfrak m \subset R$ such that
$I \subset \mathfrak m$ but $f \not\in \mathfrak m$.
Hence $\mathfrak m \in D(f) \cap V(I) = U \cap Z$ as desired.

\medskip\noindent
Conversely, suppose that $\Spec(R)$ is Jacobson.
Let $I \subset R$ be a radical ideal. Let
$J = \cap_{I \subset \mathfrak m} \mathfrak m$
be the intersection of the maximal ideals containing $I$.
Clearly $J$ is a radical ideal, $V(J) \subset V(I)$, and
$V(J)$ is the smallest closed subset of $V(I)$ containing
all the closed points of $V(I)$. By assumption we see that
$V(J) = V(I)$. But Lemma \ref{lemma-Zariski-topology}
shows there is a bijection between Zariski closed
sets and radical ideals, hence $I = J$ as desired.
\end{proof}

\end{document}
